\documentclass[10pt,a4paper]{amsart}
\usepackage[margin=19mm]{geometry}
\usepackage{amsmath,amssymb,mathtools}
\usepackage[scr=rsfs,cal=euler]{mathalfa}
\usepackage{xfrac}
\usepackage[unicode,hidelinks,bookmarks=false]{hyperref}
\newcommand{\ie}{i.e.\ }
\newcommand{\cf}{cf.\ }
\newcommand{\resp}{respectively\ }
\newcommand{\Subsection}{\subsection}
\providecommand{\nobreakdash}{\nobreak\hbox{-}}
\setlength{\emergencystretch}{2em}
\multlinegap0pt
\usepackage{amssymb}
\usepackage[all]{xy}
\usepackage{tikz}
\newtheorem{prop}{Proposition}
\newtheorem{cor}{Corollary}
\newcommand{\Q}{\mathbb{Q}}
\DeclareMathOperator{\Spec}{Spec}
\newcommand{\cO}{\mathcal{O}}
\newcommand{\km}{\mathfrak{m}}
\newcommand{\kp}{\mathfrak{p}}
\newcommand{\kq}{\mathfrak{q}}
\def\ps{\mathrm{ps}}
\newcommand{\tr}{\mathrm{tr}}
\title{Zariski density of modular points in the Eisenstein case}
\author{Xinyao Zhang}
\date{Source: arXiv:2512.21249v2 (CC BY 4.0)}
\begin{document}
\maketitle

\noindent\emph{Source and licence.} Zariski density of modular points in the Eisenstein case. Version arXiv:2512.21249v2 (CC BY 4.0), distributed by arXiv under the Creative Commons Attribution 4.0 International licence, http://creativecommons.org/licenses/by/4.0/, as recorded in the arXiv licence metadata for this exact version and shown on the abstract page https://arxiv.org/abs/2512.21249v2. Full licence notice: \texttt{licenses/arxiv-2512.21249-CC-BY-4.0.txt}.\par\smallskip

\noindent\textit{Editorial scope.} Counted unit: the article's prop dim (source0.tex lines 1110--1118). The article's complete original proof of it is delivered with it (source0.tex lines 1120--1127, one \texttt{proof} environment of 8 lines). No mathematics is written, repaired or re-derived: the statement and the proof are the article's own lines, in the article's own notation, and no retained line is substituted or re-flowed.  Retained uncounted as the source-local context the proof needs: lines 344--357; lines 469--473; lines 512--516; lines 536--538; lines 803--809; lines 1092--1094.  Nothing was omitted from any retained span, so the package carries no omission record.  No retained line contains a citation, so the wrapper carries no bibliography and no bibliography entry.  No retained line contains a figure environment, an image-inclusion command or drawing code, so no image is delivered and the package declares no asset dependency.  Editorial formatting only, in addition: an AMS \texttt{amsart} preamble that restates the article's own declarations and macros used by the retained lines, namely the environments \texttt{prop}, \texttt{cor} on separate counters, together with the standard packages those lines need.  The retained environments print in this document as Proposition 1, Corollary 1 in order of appearance, whereas the article prints the same items with the numbering of its own full document.  The complete native source of the article as distributed in the arXiv e-print for this exact version is bundled unchanged as \texttt{sources/191633/source0.tex}.
	\begin{prop}\label{com alg}
		Let $n $ be a positive integer. For any $1 \le i \le n $, assume that $R_i \in \mathcal{C}_{\mathcal{O}}$ is $\cO$-flat. (Here we say that an ideal $I$ of $R_i$ is $\cO$-flat if the ring $R_i/I$ is $\cO$-flat.)
		
		1) The completed tensor product ${\widehat{\otimes}_\cO} R_i$ is $\cO$-flat.
		
		2) If all $R_i $ are $\cO$-flat domains, $1 \le i \le n$, then $\dim {\widehat{\otimes}_\cO} R_i = 1+ \sum_{i=1}^{n} (\dim R_i -1)$.
		
		3) For any ideal $\kp_i$ of $R_i $ such that $R_i/\kp_i$ is $\cO$-flat, we have $$ 
		\widehat{\otimes}_\cO (R_i/\kp_i) =(\widehat{\otimes}_\cO R_i)/(\kp_1, ..., \kp_n). $$
		
		4) For each $i$, assume that $R_i$ is a domain of dimension $1$.  Then ${\widehat{\otimes}_\cO} R_i $ is reduced.
		
		5) Let $\{\kq_i\}$ be a set (possibly infinite) of $\cO$-flat primes of $R_1$. Then we have $(\cap \kq_i) \widehat{\otimes}_\cO R_2 = \cap (\kq_i \widehat{\otimes}_\cO R_2)$.
	\end{prop}

    \begin{prop}\label{det}
    	1) Let $\Gamma$ be the maximal pro-$p$ abelian quotient of $G_{F, \Sigma} $. Then there exists a CNL $\cO$-algebra isomorphism $R_{F}^{\ps} \to R_{F}^{\ps, \chi} \widehat{\otimes}_\cO \cO[[\Gamma]] $.
    	
    	2) Let $\chi': G_{F,\Sigma} \to \cO^\times$ be a character lifting $\bar{\chi}$. Let $R_{F}^{\ps, \chi'}$ be the universal pseudo-deformation ring parametrizing all the pseudo-deformations of $T$ with determinant $\chi'$.Then there exists a CNL $\cO$-algebra isomorphism $ R_{F}^{\ps, \chi} \to R_{F}^{\ps, \chi'}$.
    \end{prop}

    \begin{prop}\label{def of 0}
    	1) $ R_{F,0}^{\ps, \chi}$ is reduced. Each irreducible component of $ R_{F,0}^{\ps, \chi}$ is of Krull dimension at least $1+2[F:\Q]$.
    	
    	2) For any irreducible prime $\kp$ of $R_{F}^{\ps, \chi} $, the natural surjection $ R_{F}^{\ps, \chi} \twoheadrightarrow R_{F}^{\ps, \chi}/\kp$ factors through $ R_{F,0}^{\ps, \chi}$.
    \end{prop}

	\begin{cor}\label{lower}
		Assume that $\bar{\chi}$ can be extended to a character of $G_\Q$. Then we have $\dim R_{F}^{\ps, \chi} \ge 1+2[F:\Q]$. In other words, $\Spec R_{F, 0}^{\ps, \chi}$ is not empty.
	\end{cor}

    \begin{prop}\label{R_v}  Assume $\psi$ is de Rham (resp. crystalline) at $v$.
    	\begin{enumerate}
    		\item The ring $R_v$ is reduced and $\cO$-flat.
    		\item  For each minimal prime $\kq$ of $R_v$, $\kq$ is the intersection of all regular de Rham (resp. crystalline) primes containing $\kq$.
    		\item  Assume that $\tr \bar{\rho}_{\km, v} $ is not of the form $\eta + \eta\omega$ for any character $\eta$ if $p=3$. Then we have $\dim R_v \le 3$.
    	\end{enumerate}
    \end{prop}

     \begin{prop}\label{finiteness 2}
    	There exists a finite map $\widehat{\otimes}_{v \in \Sigma_p} R_{v} \to R_{F, 0}^{\ps, \chi} $.
    \end{prop}

    \begin{prop}\label{finite same dim}
    	Assume that we are not in the exceptional case.
    	
    	1) The ring $R_{F, 0}^{\ps, \chi}$ is equidimensional of dimension $1+2[F:\Q]$. In particular, $\dim R_{F}^{\ps, \chi}=1+2[F:\Q]$.
    	
    	2) Let $R_{F}^{\ps}$ be the universal pseudo-deformation ring of the pseudo-representation $1+\bar{\chi}$ for $G_{F, \Sigma}$ (without fixed determinant). Then we have $\dim R_{F}^{\ps}=2+2[F:\Q]$.
    	
    	3) For each $v|p$, we have $\dim R_v=3$.
    \end{prop}

    \begin{proof}
    	As we are not in the exceptional case, by Proposition \ref{R_v}, we know that for each $v|p$, $R_v$ is $\cO$-flat of Krull dimension at most $3$. Further, by Proposition \ref{com alg}, the ring $\widehat{\otimes}_{v \in \Sigma_p} R_{v}$ is of Krull dimension at most $1+2[F: \Q]$. On the other hand, by Proposition \ref{def of 0} and Corollary \ref{lower}, each irreducible component of $ R_{F, 0}^{\ps, \chi}$ has Krull dimension at least $1+2[F:\Q]$. Combining Proposition \ref{finiteness 2}, we have $$
    	1+2[F:\Q] \le \dim R_{F, 0}^{\ps, \chi} \le \dim \widehat{\otimes}_{v \in \Sigma_p} R_{v} \le 1+2[F:\Q].$$
    	This shows both the first and the last assertions.
    	
    	For the second one, it is direct from the first one and Proposition \ref{det} as Leopoldt's conjecture holds in our case.
    	
    \end{proof}

\end{document}
